Federated learning (FL) is a distributed machine learning paradigm that enables multiple clients to collaboratively train models without exchanging their raw data~\cite{mcmahan2017communication}, making it particularly valuable in privacy-sensitive domains such as healthcare~\cite{rieke2020future}, edge computing~\cite{lim2020federated}, and multi-institutional research~\cite{kairouz2021advances}. However, as FL deployments grow in scale and heterogeneity, they introduce significant observability challenges. A single training run may involve many clients with varying data sizes, compute capabilities, communication costs, geographic locations, and failure modes~\cite{kairouz2021advances, li2020federated}. Without a unified monitoring layer, practitioners must manually reconcile logs from servers, clients, dashboards, and experiment trackers — an error-prone process that hinders reproducibility and debugging.

In this paper, we introduce HiveWatch, a lightweight instrumentation layer for monitoring federated learning workflows. Rather than requiring a specific FL framework — such as APPFL~\cite{ryu2022appfl, li2025advances}, Flotilla~\cite{banerjee2025flotilla}, Flower~\cite{beutel2020flower}, or FedML~\cite{he2020fedml} — HiveWatch provides a framework-agnostic API that exposes common logging hooks for round starts, client updates, aggregated round results, and run completion. Our system enables users to track standard metrics including accuracy, loss, gradient norm, training time, bytes transferred, CPU and memory usage, and client location. HiveWatch achieves this through three key functions:
\begin{itemize}
    \item \textbf{Framework-Agnostic Instrumentation:} A minimal logging API that integrates with any existing FL workflow without requiring framework-specific modifications.
    \item \textbf{Per-Client Observability:} Fine-grained tracking of client-level metrics across rounds, enabling identification of stragglers, failures, and resource bottlenecks.
    \item \textbf{Unified Experiment Tracking:} A centralized view of distributed training runs, consolidating server and client logs into a single monitoring interface.
\end{itemize}

A core design goal of HiveWatch is to decouple \emph{what} gets measured from \emph{where} it is stored. Practitioners frequently move between local debugging, hosted dashboards, and self-hosted tracking servers as a project matures, and switching between these destinations should not require rewriting instrumentation code or re-deriving derived metrics such as communication volume or gradient divergence. HiveWatch addresses this by computing such derived quantities once, inside the runtime, and fanning the resulting event stream out to any combination of pluggable emitters — including Weights \& Biases, MLflow, and a local server-sent events (SSE) emitter that persists replayable JSONL and map-ready metadata artifacts. We demonstrate this design by integrating HiveWatch into three federated learning frameworks with substantially different communication stacks — APPFL~\cite{ryu2022appfl, li2025advances} (gRPC), Flower~\cite{beutel2020flower}, and NVIDIA FLARE~\cite{roth2022nvidia} (simulation-mode workflows) — using only a handful of call sites on the server and client.

The remainder of this paper is organized as follows. Section~\ref{sec:related} surveys related work in general-purpose experiment tracking and FL-specific monitoring. Section~\ref{sec:algorithm} describes HiveWatch's runtime API, emitter architecture, and map-based replay dashboard. Section~\ref{sec:eval} reports integration case studies and a microbenchmark of HiveWatch's per-call instrumentation overhead. Section~\ref{sec:conclusion} concludes and outlines future work.




% Extensive experiments on several naturally and synthetically partitioned federated datasets indicate that leveraging spot instances can reduce FL training costs by approximately 60\% compared to on-demand instances. Moreover, incorporating the proposed FedCostAware scheduling algorithm can further improve cost efficiency by achieving up to a 70\% reduction in cost, underscoring its effectiveness in optimizing resource usage.
